% Manual reference list. In-text [n] (Chapter 2) match the numbered list below.
% Author--year entries cover Chapters 3--7.

{\footnotesize\setlength{\parindent}{0pt}

\textbf{Numbered references (Chapter 2)}\par\smallskip

[1] J. S. Park, J. O'Brien, C. J. Cai, M. R. Morris, P. Liang, and M. S. Bernstein, ``Generative agents: Interactive simulacra of human behavior,'' in \textit{Proc. 36th Annual ACM Symposium on User Interface Software and Technology}, 2023, pp. 1--22.\par\smallskip
[2] J. M. Epstein and R. Axtell, \textit{Growing Artificial Societies: Social Science from the Bottom Up}. Brookings Institution Press, 1996.\par\smallskip
[3] N. Gilbert and K. Troitzsch, \textit{Simulation for the Social Scientist}. McGraw-Hill Education (UK), 2005.\par\smallskip
[4] E. Bonabeau, ``Agent-based modeling: Methods and techniques for simulating human systems,'' \textit{Proc. National Academy of Sciences}, vol. 99, no. suppl\_3, pp. 7280--7287, 2002.\par\smallskip
[5] L. P. Argyle, E. C. Busby, N. Fulda, J. R. Gubler, C. Rytting, and D. Wingate, ``Out of one, many: Using language models to simulate human samples,'' \textit{Political Analysis}, vol. 31, no. 3, pp. 337--351, 2023.\par\smallskip
[6] J. J. Horton, ``Large language models as simulated economic agents: What can we learn from homo silicus?'' National Bureau of Economic Research, Tech. Rep., 2023.\par\smallskip
[7] P. Törnberg, ``Best practices for text annotation with large language models,'' \textit{arXiv:2402.05129}, 2024.\par\smallskip
[8] M. Wooldridge, \textit{An Introduction to Multiagent Systems}. John Wiley \& Sons, 2009.\par\smallskip
[9] A. S. Rao, M. P. Georgeff et al., ``BDI agents: From theory to practice,'' in \textit{ICMAS}, vol. 95, 1995, pp. 312--319.\par\smallskip
[10] M. Wooldridge, \textit{Reasoning about Rational Agents}. The MIT Press, 2000.\par\smallskip
[11] M. E. Bratman, D. J. Israel, and M. E. Pollack, ``Plans and resource-bounded practical reasoning,'' \textit{Computational Intelligence}, vol. 4, no. 3, pp. 349--355, 1988.\par\smallskip
[12] M. P. Georgeff and A. L. Lansky, ``Reactive reasoning and planning,'' in \textit{AAAI}, vol. 87, 1987, pp. 677--682.\par\smallskip
[13] J. E. Laird, \textit{The Soar Cognitive Architecture}. MIT Press, 2019.\par\smallskip
[14] J. R. Anderson, \textit{How Can the Human Mind Occur in the Physical Universe?} Oxford University Press, 2009.\par\smallskip
[15] R. S. Sutton, A. G. Barto et al., \textit{Reinforcement Learning: An Introduction}. MIT Press Cambridge, 1998.\par\smallskip
[16] P. Hernandez-Leal, B. Kartal, and M. E. Taylor, ``A survey and critique of multiagent deep reinforcement learning,'' \textit{Autonomous Agents and Multi-Agent Systems}, vol. 33, no. 6, pp. 750--797, 2019.\par\smallskip
[17] K. Zhang, Z. Yang, and T. Başar, ``Multi-agent reinforcement learning: A selective overview of theories and algorithms,'' \textit{Handbook of Reinforcement Learning and Control}, pp. 321--384, 2021.\par\smallskip
[18] S. Russell and P. Norvig, \textit{Artificial Intelligence: A Modern Approach, Global Edition}, 4th ed. Pearson Education, 2021.\par\smallskip
[19] A. Vaswani, N. Shazeer, N. Parmar, J. Uszkoreit, L. Jones, A. N. Gomez, Ł. Kaiser, and I. Polosukhin, ``Attention is all you need,'' \textit{Advances in Neural Information Processing Systems}, vol. 30, 2017.\par\smallskip
[20] S. Bubeck, V. Chandrasekaran, R. Eldan, J. Gehrke, E. Horvitz, E. Kamar, P. Lee, Y. T. Lee, Y. Li, S. Lundberg et al., ``Sparks of artificial general intelligence: Early experiments with GPT-4,'' \textit{arXiv:2303.12712}, 2023.\par\smallskip
[21] J. Wei, X. Wang, D. Schuurmans, M. Bosma, F. Xia, E. Chi, Q. V. Le, D. Zhou et al., ``Chain-of-thought prompting elicits reasoning in large language models,'' \textit{Advances in Neural Information Processing Systems}, vol. 35, pp. 24824--24837, 2022.\par\smallskip
[22] P. Lewis, E. Perez, A. Piktus, F. Petroni, V. Karpukhin, N. Goyal, H. Küttler, M. Lewis, W.-t. Yih, T. Rocktäschel et al., ``Retrieval-augmented generation for knowledge-intensive NLP tasks,'' \textit{Advances in Neural Information Processing Systems}, vol. 33, pp. 9459--9474, 2020.\par\smallskip
[23] S. Yao, J. Zhao, D. Yu, N. Du, I. Shafran, K. R. Narasimhan, and Y. Cao, ``ReAct: Synergizing reasoning and acting in language models,'' in \textit{The Eleventh International Conference on Learning Representations}, 2022.\par\smallskip
[24] S. Yao, J. Zhao, D. Yu, N. Du, I. Shafran, K. Narasimhan, and Y. Cao, ``ReAct: Synergizing reasoning and acting in language models,'' 2022, project website. Available: https://react-lm.github.io/\par\smallskip
[25] R. J. Beckman, K. A. Baggerly, and M. D. McKay, ``Creating synthetic baseline populations,'' \textit{Transportation Research Part A: Policy and Practice}, vol. 30, no. 6, pp. 415--429, 1996.\par\smallskip
[26] D. J. Watts and S. H. Strogatz, ``Collective dynamics of `small-world' networks,'' \textit{Nature}, vol. 393, pp. 440--442, 1998.\par\smallskip
[27] J. S. Park, C. Q. Zou, A. Shaw, B. M. Hill, C. J. Cai, M. R. Morris, R. Willer, P. Liang, and M. S. Bernstein, ``Generative agent simulations of 1,000 people,'' \textit{arXiv:2411.10109}, 2024.\par\smallskip
[28] G. V. Aher, R. I. Arriaga, and A. T. Kalai, ``Using large language models to simulate multiple humans and replicate human subject studies,'' in \textit{Proc. 40th International Conference on Machine Learning (ICML)}, 2023.\par\smallskip
[29] K. Li, T. Liu et al., ``Measuring and controlling persona drift in language model dialogues,'' \textit{arXiv:2402.10962}, 2024.\par\smallskip
[30] J. Bisbee, J. D. Clinton, C. Dorff, B. Kenkel, and J. M. Larson, ``Synthetic replacements for human survey data? The perils of large language models,'' \textit{Political Analysis}, vol. 32, 2024.\par\smallskip
[31] C. Gao, X. Lan, N. Li, Y. Yuan, J. Ding, Z. Zhou, F. Xu, and Y. Li, ``Large language models empowered agent-based modeling and simulation: A survey and perspectives,'' \textit{arXiv:2312.11970}, 2024.\par\smallskip
[32] A. J. Collins, P. Sokolowski, and C. Banks, ``Towards validation and verification of agent-based and complex social systems models,'' \textit{Journal of Defense Modeling and Simulation}, 2015.\par\smallskip
[33] D. J. C. MacKay, \textit{Information Theory, Inference, and Learning Algorithms}. Cambridge University Press, 2003.\par\smallskip
[34] M. Arjovsky, S. Chintala, and L. Bottou, ``Wasserstein generative adversarial networks,'' in \textit{Proc. 34th International Conference on Machine Learning (ICML)}, 2017, pp. 214--223.\par\smallskip
[35] S. Manzan, ``Wasserstein distance and the distributional analysis of economic data,'' working paper, 2021.\par\smallskip
[36] F. A. Cowell, \textit{Measuring Inequality}, 3rd ed. Oxford University Press, 2011.\par\medskip

\textbf{Author--year references (Chapters 3--7)}\par\smallskip

Acemoglu, D. \& Robinson, J.A. (2012). \textit{Why Nations Fail}. Crown Publishers.\par\smallskip
Aher, G. et al. (2023). ``Using Large Language Models to Simulate Multiple Humans and Replicate Human Subject Studies.'' \textit{ICML 2023}.\par\smallskip
Argyle, L.P. et al. (2023). ``Out of One, Many: Using Language Models to Simulate Human Samples.'' \textit{Political Analysis}, 31(3).\par\smallskip
Axelrod, R. (1984). \textit{The Evolution of Cooperation}. Basic Books.\par\smallskip
Axelrod, R. (1997). \textit{The Complexity of Cooperation}. Princeton University Press.\par\smallskip
Barabási, A.-L. \& Albert, R. (1999). ``Emergence of Scaling in Random Networks.'' \textit{Science}, 286(5439), 509--512.\par\smallskip
Berg, J., Dickhaut, J. \& McCabe, K. (1995). ``Trust, Reciprocity, and Social History.'' \textit{Games and Economic Behavior}, 10(1), 122--142.\par\smallskip
Chaudhuri, A. (2011). ``Sustaining cooperation in laboratory public goods experiments: a selective survey of the literature.'' \textit{Experimental Economics}, 14(1), 47--83.\par\smallskip
Clauset, A., Shalizi, C.R. \& Newman, M.E.J. (2009). ``Power-Law Distributions in Empirical Data.'' \textit{SIAM Review}, 51(4), 661--703.\par\smallskip
Cohen, J. (1988). \textit{Statistical Power Analysis for the Behavioral Sciences} (2nd ed.). Lawrence Erlbaum Associates.\par\smallskip
Cover, T.M. \& Thomas, J.A. (2006). \textit{Elements of Information Theory} (2nd ed.). Wiley.\par\smallskip
Epstein, J.M. (2006). \textit{Generative Social Science}. Princeton University Press.\par\smallskip
Epstein, J.M. \& Axtell, R. (1996). \textit{Growing Artificial Societies: Social Science from the Bottom Up}. MIT Press.\par\smallskip
European Social Survey ERIC (2024). \textit{ESS Round 11 -- 2023}. Data file edition.\par\smallskip
Falk, A. et al. (2018). ``Global Evidence on Economic Preferences.'' \textit{Quarterly Journal of Economics}, 133(4), 1645--1692.\par\smallskip
Fehr, E. \& Gächter, S. (2000). ``Cooperation and Punishment in Public Goods Experiments.'' \textit{American Economic Review}, 90(4), 980--994.\par\smallskip
Fontana, M., Pierri, F. \& Aiello, L.M. (2024). ``Nicer Than Humans: How do Large Language Models Behave in the Prisoner's Dilemma?'' \textit{arXiv:2406.13605}.\par\smallskip
Gao, C. et al. (2023). ``S3: Social-Network Simulation System with Large Language Model-Empowered Agents.'' \textit{arXiv:2307.14984}.\par\smallskip
Gibbs, A.L. \& Su, F.E. (2002). ``On Choosing and Bounding Probability Metrics.'' \textit{International Statistical Review}, 70(3), 419--435.\par\smallskip
Glaeser, E.L. et al. (2000). ``Measuring Trust.'' \textit{Quarterly Journal of Economics}, 115(3), 811--846.\par\smallskip
Hedges, L.V. (1981). ``Distribution Theory for Glass's Estimator of Effect Size and Related Estimators.'' \textit{Journal of Educational Statistics}, 6(2), 107--128.\par\smallskip
Henrich, J. et al. (2010). ``Markets, Religion, Community Size, and the Evolution of Fairness and Punishment.'' \textit{Science}, 327(5972), 1480--1484.\par\smallskip
Hernán, M.A. \& Robins, J.M. (2020). \textit{Causal Inference: What If}. Chapman \& Hall/CRC.\par\smallskip
Herrmann, B., Thöni, C. \& Gächter, S. (2008). ``Antisocial Punishment Across Societies.'' \textit{Science}, 319(5868), 1362--1367.\par\smallskip
Holland, J.H. (1992). \textit{Adaptation in Natural and Artificial Systems}. MIT Press.\par\smallskip
Horton, J.J. (2023). ``Large Language Models as Simulated Economic Agents: What Can We Learn from Homo Silicus?'' \textit{NBER Working Paper 31122}.\par\smallskip
Inglehart, R. \& Welzel, C. (2010). ``Changing Mass Priorities: The Link between Modernization and Democracy.'' \textit{Perspectives on Politics}, 8(2), 551--567.\par\smallskip
Kauffman, S. (1993). \textit{The Origins of Order}. Oxford University Press.\par\smallskip
Lewis, P. et al. (2020). ``Retrieval-Augmented Generation for Knowledge-Intensive NLP Tasks.'' \textit{NeurIPS 2020}.\par\smallskip
Li, G. et al. (2023). ``CAMEL: Communicative Agents for `Mind' Exploration of Large Language Model Society.'' \textit{NeurIPS 2023}. arXiv:2303.17760.\par\smallskip
Liu, X. et al. (2024). ``AgentBench: Evaluating LLMs as Agents.'' \textit{ICLR 2024}. arXiv:2308.03688.\par\smallskip
Manning, J.P. et al. (2024). ``Automated Social Science: Language Models as Scientist and Subjects.'' \textit{arXiv:2404.11794}.\par\smallskip
Mou, X. et al. (2024). ``Individual and Collective Behavior Simulation of Large Language Model Agents.'' \textit{arXiv:2402.16871}.\par\smallskip
Münker, L., Schwager, I. \& Rettinger, A. (2025). ``Don't Trust Generative Agents to Mimic Communication on Social Networks Unless You Benchmarked their Empirical Realism.'' \textit{arXiv:2506.21974}.\par\smallskip
Nowak, M.A. \& May, R.M. (1992). ``Evolutionary Games and Spatial Chaos.'' \textit{Nature}, 359, 826--829.\par\smallskip
Ouyang, L. et al. (2022). ``Training Language Models to Follow Instructions with Human Feedback.'' \textit{NeurIPS 2022}.\par\smallskip
Park, J.S. et al. (2023). ``Generative Agents: Interactive Simulacra of Human Behavior.'' \textit{UIST 2023}.\par\smallskip
Piketty, T. (2014). \textit{Capital in the Twenty-First Century}. Harvard University Press.\par\smallskip
Rossetti, G. et al. (2024). ``Y Social: An LLM-Powered Social Media Digital Twin.'' \textit{arXiv:2408.00818}.\par\smallskip
Schelling, T.C. (1971). ``Dynamic Models of Segregation.'' \textit{Journal of Mathematical Sociology}, 1(2), 143--186.\par\smallskip
Sharma, M. et al. (2023). ``Towards Understanding Sycophancy in Language Models.'' \textit{arXiv:2310.13548}.\par\smallskip
Tu, T. et al. (2024). ``From Single Agent to Multi-Agent: Exploring the Landscape of LLM Agent Society.'' \textit{arXiv:2402.01659}.\par\smallskip
VanderWeele, T.J. \& Ding, P. (2017). ``Sensitivity Analysis in Observational Research: Introducing the E-value.'' \textit{Annals of Internal Medicine}, 167(4), 268--274.\par\smallskip
Wang, L. et al. (2024). ``A Survey on Large Language Model-based Autonomous Agents.'' \textit{Frontiers of Computer Science}, 18(6).\par\smallskip
Watts, D.J. \& Strogatz, S.H. (1998). ``Collective Dynamics of `Small-World' Networks.'' \textit{Nature}, 393, 440--442.\par\smallskip
Zheng, J. et al. (2024). ``ChatGPT is a Knowledgeable but Inexperienced Solver.'' \textit{NAACL 2024}. arXiv:2303.16421.\par\smallskip

}
